\chapter{Related Work}\label{chap:related-work}

\setlength{\parindent}{1em}

\subsection{Formal Theorem Proving and Proof Assistants}
\label{sec:rw_formal}

Proof assistants such as Lean~4 \cite{LeanTheoremProver}, Rocq \cite{RocqTheoremProver}, and Isabelle \cite{IsabelleTheoremProver} have transformed mathematical practice by making the correctness of proofs explicit and mechanized. Community libraries such as Lean's Mathlib \cite{Mathlib} now grow at a pace driven by both increasing human contributions and by recent advances in neural theorem proving and autoformalization. This rapid expansion stresses the maintainability, coherence, and long-term usability of libraries, due to proofs that are often heterogeneous in style and clarity. Moreover, the growth rate of formal libraries already exceeds what human reviewers can reliably curate --- a discrepancy expected to widen due to increasing volumes of machine-generated proofs, which, even when formally correct, do not carry similar guarantees as to their quality, modularity, or understandability \cite{SeedProver}.

There is also a rich literature on the varied styles of human formal proofs (e.g., \cite{serge2010,Wiedijk}), noting that proofs written for different purposes --- readability, brevity, machine-aided search --- take structurally different forms. At a high level, proof assistants constitute a sound verifier in a prover-verifier game \cite{prover-verifier-games}, suggesting that a machine-learning-based prover that interfaces with such a verifier is a natural next step for formal reasoning systems.

\subsection{Neural Theorem Proving}
\label{sec:rw_ntp}

A typical approach to developing machine learning provers \cite{polu2020generative} is to train on a large corpus of mathematical proofs such as Mathlib \cite{han2022proof,polu2022formal,lample2022hypertree,yang2023leandojo}. A model learns from the distribution of proofs in the corpus, such as Mathlib-style proofs. Recently, the AlphaProof \cite{alphaproof} system demonstrated solving IMO problems by producing proofs with an arcane, non-human structure and syntax, and subsequent systems have continued to push this frontier \cite{achim2025aristotleimolevelautomatedtheorem, SeedProver}.

Many systems additionally utilize neurosymbolic augmentation, providing a generative prover model with information gathered from within the proof environment \cite{leandojo, ImProver, GodelProver}. However, both formal and informal proofs generated by current LLM-based systems often suffer from stylistic irregularities even when sound, including redundant steps or structures that do not clearly represent the broader logical argument \cite{frieder2025datamathematicalcopilotsbetter}.

The miniCTX benchmark \cite{hu2024minictxneuraltheoremproving} evaluates neural provers on theorems drawn from context-rich real-world Lean files, providing a more realistic evaluation setting than miniF2F \cite{zheng2022miniff}. ImProver~2 uses the miniCTX-v2 extension of this benchmark as its primary test set, enabling evaluation on genuine research-level mathematics.

\subsection{Proof Optimization}
\label{sec:rw_optimization}

Automated proof optimization --- the task of rewriting a verified proof to make it better with respect to some criterion while preserving correctness --- was introduced as a formal task by ImProver \cite{ImProver} (Chapter~\ref{chap:improver} of this thesis). ImProver created a system capable of optimizing towards multiple metrics of improvement; however, it relied on general-purpose closed-source models (GPT-4o), leading to substantial deployment costs and limited ability to improve performance beyond these models' baseline.

Concurrent work by \cite{gu2025proofoptimizer} focused solely on optimizing the token count of proofs according to a complex tokenizer intended to reduce compilation time. This work does not examine other metrics, leaving out important use cases and limiting its utility to research mathematicians.

Low library quality also decreases its utility as training data: modern theorem provers and autoformalizers increasingly train on these very corpora, so the structure and readability of proofs directly shape downstream prover performance \cite{gu2025proofoptimizer}.

ImProver~2 \cite{ImProver2} (Chapter~\ref{chap:improver2} of this thesis) addresses the cost and scalability limitations of ImProver by training a small, specialized 7B model that matches or exceeds frontier models on structural metrics.

\subsection{LLM Agents and Prompting}
\label{sec:rw_agents}

ImProver draws on a range of techniques from the LLM agent and prompting literature. Chain-of-thought prompting \cite{wei2022chain} makes intermediate reasoning steps explicit, improving performance on complex tasks; ImProver's Chain-of-States prompting (\S\ref{ssec:cos_imp}) is the formal-proof analogue. Full proof generation from sketches \cite{jiang2023draft,first2023baldur} and conditioning on example proofs \cite{jiang2023draft} form the basis of ImProver's retrieval and prompting strategy. Preceding file context \cite{first2023baldur,hu2024minictxneuraltheoremproving} improves proof generation quality by providing relevant library context.

\subsection{Retrieval-Augmented Generation}
\label{sec:rw_rag}

ImProver incorporates retrieval-augmented generation (RAG) from both code generation and theorem proving. Documentation retrieval for code generation \cite{zhou2023docprompting} provides syntactic guidance via the TPiL handbook. Retrieval from the Mathlib library \cite{yang2023leandojo,thakur2024incontext} enables the model to find relevant lemmas to use in its rewrites. ImProver uses Maximum Marginal Relevance \cite{10.1145/290941.291025} to select diverse, relevant examples and documents. ImProver~2 replaces explicit retrieval with context extraction from the proof environment (\S\ref{sec:context_extraction}), which is more targeted and requires no separate retrieval index.

\subsection{Preference Optimization and Reinforcement Learning}
\label{sec:rw_rl}

ImProver~2 builds upon a substantial body of work on preference optimization and reinforcement learning from feedback. Direct Preference Optimization (DPO) \cite{DPO} provides a stable, offline alternative to RLHF by framing preference learning as a classification problem. Iterative Reasoning Preference Optimization (IRPO) \cite{IRPO} extends DPO with an NLL regularization term that improves training stability and sample efficiency.

Expert iteration \cite{ExpertIteration} bootstraps model capabilities by iteratively generating training data and retraining, enabling self-improvement without external supervision. Group Relative Policy Optimization (GRPO) \cite{GRPO} is an alternative RL policy that learns from groups of candidate solutions; ImProver~2 uses pairwise preference data instead, which provides a denser reward signal for the limited amount of publicly available Lean training data.

Model collapse \cite{modelCollapse} is a failure mode in which indiscriminate consumption of self-generated training data causes distribution collapse. ImProver~2's replay buffer is specifically designed to prevent this by filtering and combining new and old data across training iterations.

\subsection{Neurosymbolic Methods}
\label{sec:rw_neurosymbolic}

Neurosymbolic methods combine neural generation with formal symbolic reasoning, leveraging the strengths of both. In the theorem proving context, this typically involves using formal proof checkers to validate generated proofs and extracting symbolic information from the proof environment to guide generation.

ImProver~2 extends neurosymbolic augmentation in three directions. Context extraction (\S\ref{sec:context_extraction}) collects the signatures and documentation of directly referenced lemmas and definitions to inform generation. Chain-of-States prompting (\S\ref{ssec:cos_imp}) provides goal-state traces extracted from Lean's InfoTree structures. Auto-informalization (\S\ref{sec:generation}) uses natural language translation to provide a higher-level representation of the proof, inspired by the draft-sketch-prove approach \cite{DraftSketchProve} and work on auto-formalization \cite{autoinformalization}.

The combination of these three augmentation sources consistently improves performance across both small and frontier models, suggesting that formal structure exposure is complementary to raw model capability.
